% This document was generated by an LLM.

\documentclass{self-check}

\title{Self-Check Questions: Absolute Values in Differential and Integral Calculus}
\author{}
\date{}

\begin{document}
\maketitle

\noindent\textit{Instructions.} Answer without consulting notes. Use formulas involving absolute values or the sign function rather than rewriting functions piecewise. Use difference quotients at points where differentiation rules do not apply. Justify all formulas and state domains, points of non-differentiability, and constants of integration. Where an integral is indefinite, give an antiderivative on each connected interval of its domain. No solutions are included.

\section*{1. Definitions and elementary identities}
\begin{enumerate}
\item State the piecewise definition of $|x|$. Prove that $|x|^2=x^2$ and $|x|=\sqrt{x^2}$.
\item Prove that $x/|x|=|x|/x$ for $x\ne0$. Why are neither side defined at zero?
\item Define the sign function $\operatorname{sgn}(x)$. Express it using absolute values for $x\ne0$, and explain how its value at zero is conventionally defined.
\item Prove that $|x|=x\operatorname{sgn}(x)$ for all real $x$ (using $\operatorname{sgn}(0)=0$).
\end{enumerate}

\section*{2. Derivatives: fundamentals}
\begin{enumerate}[resume]
\item Find $\dfrac{d}{dx}|x|$ for $x\ne0$ from the definition of the derivative. Express the result using absolute values or the sign function.
\item Prove from the one-sided difference quotients that $|x|$ is not differentiable at $x=0$.
\end{enumerate}

\section*{3. Chain rule and points of non-differentiability}
\begin{enumerate}[resume]
\item Use the chain rule and $|x|=\sqrt{x^2}$ to recover the derivative of $|x|$ for $x\ne0$. Why does this calculation not establish differentiability at zero?
\item Suppose $g$ is differentiable at $x$ and $g(x)\ne0$. Prove that $\dfrac{d}{dx}|g(x)|=\dfrac{g(x)}{|g(x)|}g'(x)$.
\item Find $\dfrac{d}{dx}\ln|x|$, including its domain. Prove your result using the chain rule.
\item Find $\dfrac{d}{dx}\ln|x^2-4|$ and identify all excluded points.
\end{enumerate}

\subsection*{Differentiability at zeros}
Unless a question explicitly weakens the assumptions, let $g$ be a real-valued function defined on an open interval $I$ containing $c$, differentiable at $c$, with $g(c)=0$. Set $f(x)=|g(x)|$. Write $f'_+(c)$ and $f'_-(c)$ for the right-hand and left-hand derivatives, respectively.
Here, $g$ \emph{changes sign across $c$} if, for some $\delta>0$, it is strictly positive on $(c-\delta,c)$ and strictly negative on $(c,c+\delta)$, or vice versa, with both intervals contained in $I$.

\begin{enumerate}[resume]
\item Write the difference quotient defining $f'(c)$ and simplify it using $g(c)=0$. State separately the limits defining $f'_+(c)$ and $f'_-(c)$.
\item What sign must this difference quotient have for $h>0$? For $h<0$? Prove that if $f'(c)$ exists, it must equal zero. Does this prove that $f'(c)$ exists?
\item Explain why $c$ is a global minimum point of $f$ on $I$. State Fermat's theorem with its hypotheses and apply it conditionally at $c$. What necessary condition does it give for $f$ to be differentiable there, and why does this argument alone not prove differentiability?
\item For each of $g(x)=x$, $g(x)=x^2$, and $g(x)=x^3$ at $c=0$, determine $g'(0)$, whether $g$ changes sign across zero, and whether $|g|$ is differentiable there. Justify the behavior of $|g|$ at zero from difference quotients.
\item Consider the assertions that a sign change of $g$ across $c$ is necessary for $|g|$ to be non-differentiable at $c$, and that it is sufficient. Which assertion can the preceding examples disprove? Which assertion still requires a proof or another counterexample?
\item Rewrite the difference quotient for $f$ when $h>0$ in terms of the difference quotient for $g$. Use differentiability of $g$ at $c$ and continuity of the absolute value function to prove that $f'_+(c)$ exists and express it in terms of $g'(c)$.
\item Carry out the corresponding argument for $h<0$. Prove that $f'_-(c)$ exists and express it in terms of $g'(c)$. Account carefully for the sign of the denominator.
\item State and prove a necessary and sufficient condition on $g'(c)$ for $|g|$ to be differentiable at $c$, including all hypotheses and the value of the derivative when it exists. Prove both implications using the one-sided derivatives.
\item In the case $g'(c)=0$, give a direct proof from the two-sided difference quotient that $|g|$ is differentiable at $c$. Is continuity of $g'$ near $c$, or even existence of $g'$ away from $c$, needed? Justify your answer.
\item Suppose $g'(c)\ne0$. Prove that there is a $\delta>0$ such that $g$ has a strict constant sign on $(c-\delta,c)$ and the opposite strict sign on $(c,c+\delta)$, with both intervals contained in $I$. Determine which sign occurs on each side according to the sign of $g'(c)$.
\item Suppose instead that for some $\delta>0$ with $(c-\delta,c+\delta)\subset I$, either $g(x)\ge0$ for every $x\in(c-\delta,c+\delta)$ or $g(x)\le0$ throughout that interval. Prove that $g'(c)=0$. What follows about differentiability of $|g|$ at $c$?
\item Return to the two sign-change assertions. Under the standing differentiability assumptions, decide whether a sign change across $c$ is necessary, sufficient, both, or neither for non-differentiability of $|g|$. Prove every valid implication and give a counterexample to every invalid one. Distinguish changing sign from merely having a zero.
\item For each positive integer $n$, find the derivative of $|x^n|$ for $x\ne0$ and decide whether it is differentiable at zero. How do the parity of $n$ and the value of $n$ affect sign changes of $x^n$ and differentiability of its absolute value?
\item Find $\dfrac{d}{dx}|x^2-1|$ on every interval where it exists. At each zero of $x^2-1$, determine the two one-sided derivatives of its absolute value and decide whether the two-sided derivative exists.
\item Determine all zeros of $\sin x$ and analyze differentiability of $|\sin x|$ at each one using the theorem you proved. Find the derivative elsewhere and justify that you have identified every point of non-differentiability.
\item Now assume only that $g$ is continuous at $c$, with $g(c)=0$.
\begin{enumerate}[label=(\alph*)]
\item Does continuity alone ensure that $|g|$ is differentiable at $c$? Prove this or give a counterexample.
\item If $|g|$ is differentiable at $c$, must $g$ be differentiable there? What can be concluded about $g'(c)$? Give a proof or a counterexample.
\item Under these weaker assumptions, is a sign change necessary for non-differentiability of $|g|$? Is it sufficient? Prove each claim you accept and give a counterexample to each claim you reject.
\end{enumerate}
\item Finally, replace differentiability of $g$ at $c$ by the existence of finite one-sided derivatives $g'_+(c)=a$ and $g'_-(c)=b$, retaining $g(c)=0$ and the open interval domain.
\begin{enumerate}[label=(\alph*)]
\item Derive formulas for $f'_+(c)$ and $f'_-(c)$ in terms of $a$ and $b$.
\item Find and prove a necessary and sufficient condition on $a$ and $b$ for $|g|$ to be differentiable at $c$.
\item Can $|g|$ be differentiable at $c$ when $a\ne b$? Prove your answer or give a counterexample, and explain its relation to the original theorem.
\end{enumerate}
\end{enumerate}

\subsection*{Further applications}
\begin{enumerate}[resume]
\item Find $\dfrac{d}{dx}|x-a|$ and specify where the derivative exists, for a fixed $a\in\mathbb R$.
\item Find $\dfrac{d}{dx}|3x-6|$ and describe the derivative at $x=2$.
\item Find $\dfrac{d}{dx}|x|^3$ and determine whether the derivative exists at zero.
\item Find $\dfrac{d}{dx}x|x|$ and determine whether the derivative exists at zero.
\item Differentiate $|x|+|x-1|$ wherever possible. Identify intervals where the function is constant.
\item Differentiate $\sqrt{|x|}$ for $x\ne0$. Investigate the one-sided difference quotients at zero.
\item Differentiate $\bigl||x|-1\bigr|$ using the chain rule wherever it applies, and identify all corners using one-sided derivatives.
\item For $a\ne0$, determine whether $|x|^a$ is differentiable at zero when it is defined there. Examine separately $a>1$, $a=1$, $0<a<1$, and $a<0$.
\end{enumerate}

\section*{4. Indefinite integrals}
\begin{enumerate}[resume]
\item Find $\displaystyle\int |x|\,dx$ by reversing the differentiation of $x|x|$. Give a global antiderivative.
\item Verify your preceding answer by differentiation, including at $x=0$.
\item Find $\displaystyle\int |x-a|\,dx$ for a real constant $a$.
\item Find $\displaystyle\int |2x-3|\,dx$ using substitution and your formula for $\int |x|\,dx$.
\item Find $\displaystyle\int x|x|\,dx$ and check the answer by differentiation.
\item Find $\displaystyle\int |x|^3\,dx$.
\item Find $\displaystyle\int |x^2-1|\,dx$. Express a global antiderivative using absolute values or sign functions, and verify differentiability at $x=\pm1$.
\item Find $\displaystyle\int (|x|+|x-1|)\,dx$ and construct a single global antiderivative.
\item Find $\displaystyle\int\frac{1}{x}\,dx$ on $x>0$ and $x<0$. Explain precisely the meaning of $\ln|x|+C$ on the disconnected domain $\mathbb R\setminus\{0\}$.
\item Find $\displaystyle\int\frac{1}{|x|}\,dx$ on each connected component of its domain. Is $\ln|x|+C$ an answer?
\item Find $\displaystyle\int\frac{x}{|x|}\,dx$ on $x\ne0$. Explain the relation to $|x|$.
\item Find $\displaystyle\int\frac{2x}{x^2-1}\,dx$ and give the full real domain of the integrand.
\item Find $\displaystyle\int\frac{g'(x)}{g(x)}\,dx$ assuming $g$ is differentiable and nonzero on an interval. State any regularity assumptions needed for your reasoning.
\end{enumerate}

\section*{5. Definite integrals and interpretation}
\begin{enumerate}[resume]
\item Compute $\displaystyle\int_{-2}^{3}|x|\,dx$ by geometry and by the Fundamental Theorem of Calculus.
\item Compute $\displaystyle\int_{-2}^{2}|x^2-1|\,dx$ using the global antiderivative from the preceding section.
\item Compute $\displaystyle\int_{-1}^{2}|2x-1|\,dx$.
\item Prove that $\displaystyle\int_{-a}^{a}|x|\,dx=a^2$ for $a>0$.
\item For continuous $f$, prove $\displaystyle\left|\int_a^b f(x)\,dx\right|\le\int_a^b|f(x)|\,dx$ when $a<b$. What does equality require?
\item Explain the distinction between $\displaystyle\int_a^b|f(x)|\,dx$ and $\displaystyle\left|\int_a^b f(x)\,dx\right|$, using a concrete example where they differ.
\end{enumerate}

\section*{6. Synthesis and proof challenges}
\begin{enumerate}[resume]
\item Prove that $F(x)=\int_0^x|t|\,dt$ is differentiable at every real $x$, and find $F'(x)$.
\item Prove that $f(x)=|x|$ is continuous everywhere despite being non-differentiable at zero.
\item Investigate $f(x)=|x|^p$ at zero for $p>0$: for which values of $p$ does $f'(0)$ exist? For which values does $f''(0)$ exist?
\item Is $\displaystyle\frac{d}{dx}|g(x)|=\left|g'(x)\right|$ generally true? Give a counterexample and state the correct formula where $g(x)\ne0$.
\item Is $\displaystyle\int|f(x)|\,dx=\left|\int f(x)\,dx\right|+C$ generally true? Give a counterexample and explain the role of additive constants.
\end{enumerate}

\end{document}
